\section{Relative degeneration and the sign of a handle}\label{app:spinproof}

We prove Theorem~\ref{thm:spin}. Deformation of the metric and the
normal rescaling theorem in \cite[Proposition~15.4 and Appendices~D--E]{WueUnramified} identify the
fixed-target integral for any of the stated germs with its value at a
harmonic linear representative. At each fixed degree and source genus
this continuation takes place over the compact stable-map stack. The
vertical area identity on positive normalizations, followed by
matching on contracted components, keeps the normal family injective.

\subsection{The orientable normal bundle}
In the splitting \(E=L_+\oplus L_-\), the two degrees are
\(b_0+s\) and \(b_0-s\). Choose an orthogonal product metric.
Let \(\pi\) and \(f\) denote the universal source and universal map
over \(\overline{\mathcal M}_k(C,m)\), and put \(W_\pm=-R\pi_*f^*L_\pm\),
\(n=k-1-mb_0\). The antilinear term of \(D\) exchanges the lines;
the family from Dolbeault to the positive operator is complex-linear
in the complex structure \((i,-i)\). At its positive end the kernel is
zero. At its Dolbeault end the orientation change on the second
summand is \((-1)^{n+ms}\), since this is the parity of its
Dolbeault index. Stabilization by forms supported away from the nodes
continues the family index, including the complex equations at the
nodes, across jumps in the individual cohomology bundles. Thus
\begin{equation}\label{eq:oriented}
 e(V_D)=(-1)^{n+ms}c_{2n}(W_++\overline W_-).
\end{equation}
The right side is the cohomological degree \(2n\) term (in complex
degrees) of the equivariant Euler product for weights \(+1,-1\).
All other terms integrate to zero. Bryan--Pandharipande's
antidiagonal formula \cite[Corollary~7.2]{bp08} now gives the first
line of the theorem. Exchanging the weight convention transposes all
diagrams and preserves the sum.

\subsection{Two smoothings of the spin pair}
Pinch a nonseparating circle on \(C\); denote the normalization of
the central target by \(B\) and its ordered inverse images of the
node by \(p,q\). For each \(\epsilon\) we need one smoothing of a
pair of spins whose ratio has nontrivial holonomy around the loop
of the dual graph. The unmarked component classification in
\cite[Appendix~E]{WueUnramified} then applies to every pair of that sign.
For \(\epsilon=1\) use the exceptional spin at an index-two node.
The two spin structures restrict to the same theta characteristic
on \(B\) and differ by a sign in sewing at \(q\). Over an exceptional
rational component the spin has degree one, so its sections interpolate
the two end values uniquely. Both extensions thus have the same
parity. For \(\epsilon=-1\), take the two locally free extensions of
\(A^2=K_B(p+q)\) at an ordinary node. Riemann--Roch gives
\(h^0(A)-h^0(A(-p-q))=1\); the line of end values is isotropic for
the sum of squared residues, by the residue theorem. Hence exactly one
of the two isotropic clutch lines admits the extra section. The two
parities differ. The roots smooth after a local base change, by
smoothness of the relative Picard space at a line and étaleness of
squaring; parity is preserved in spin smoothing \cite[\S6]{cornalba}. In both cases their
ratio has trivial inertia and its central double torsor is connected
by the nontrivial clutch. This persists under smoothing. See
Calculation~D1 of \cite{WueSupplement} for the square frames.
Use \(d=2\) for the exceptional spin and \(d=1\) for the locally free
spin.

After \(\delta=z^2\), resolve the local smoothing \(UV=\delta\)
before taking the stabilizer quotient. The two central components
are the \(d\)-th root of \(B\) at \(p,q\) and the two-ended
rational component \([\PP^1/\mu_d]\) with its rotation action.
Both have coarse covering degree \(m\), by intersection with the
exceptional component in the surface. We use the simultaneous
relative expansions at \(p,q\) of \cite[\S\S4.4--4.7]{af}.
In a frame whose square is the logarithmic differential, take the
antilinear term proportional to \(U\,d\bar U-V\,d\bar V\) on the neck.
Its coefficient relative to the conjugate logarithmic differential is
\(|U|^2+|V|^2\) on a smooth fiber and can be joined by positive
interpolation to the metric off the neck. Pullback by contraction
makes it zero on the inserted components.

Stokes' identity on a main normalization gives nonnegative interior
energy and opposite end residues weighted by the coarse contact
orders; orbifold integration cancels the uniformizer index. The
square of a holomorphic solution on an inserted component is
constant. Conservation of contact order thus cancels the end
residues on the glued source. The solution vanishes on each
nonconstant main component, then on all components by matching.
In particular contracted subcurves of arbitrarily high genus do not
create a kernel. Stabilization with forms off the nodes gives an
oriented cokernel bundle on the plumbing charts. The antilinear term
is the pullback of a bounded ambient coefficient on each half-neck,
so the nodal estimates of \cite[Appendix~B]{bai} apply with this term.
The family Euler product restricts under virtual Gysin to the paired
Euler insertion on the central splitting. The expanded family and
virtual restriction are described in
Calculations~D6--D7 of \cite{WueSupplement}.

\subsection{Eliminating the two-ended component}
On a connected exceptional relative source of arithmetic genus \(j\)
let \(S\) denote the pulled-back sign line. Its space of constants
is a line or zero according as its double torsor is trivial or
connected; for \(d=1\) it is always a line. These alternatives are
locally constant on the relative stack. The normalization sequence
gives the two real cokernels a common complex quotient
\(A=R^1\pi_*S\). Mumford's Hodge relation on the coarse double torsor,
divided by the same relation on its invariant summand, gives
\(c(A)c(A^\vee)=1\). In positive rank \(c_{\rm top}(A)^2=0\), so the
paired insertion vanishes. In rank zero eliminate the surjective
exceptional Dolbeault equations. The remaining equations involve the
main operator, its evaluations at the ends and one constant for each
trivial-torsor vertex. Pass to the finite étale cover framing the
invariant end lines by square-compatible frames. The residual choices
are finite signs, and on every such sector the insertion is pulled
back from the main relative space. Finite transfer thus reduces the
exceptional integral to its virtual class paired against \(1\).

The log tangent of the two-ended target is trivial. A connected
positive-degree exceptional source with \(a\) ends meets both colors,
so its virtual dimension \(2j-2+a\) is nonnegative. Integration against
\(1\) vanishes in positive dimension. Only genus zero with two ends
survives. Stability and Riemann--Hurwitz then leave the power map
of degree \(k\), branched only at the two ends; unbranched rubber
levels give the same relative map. The two profiles agree, say
\(\mu=(k_1,\ldots,k_r)\). The following multiplicity allows repeated
parts as well as even contacts.
\begin{lemma}\label{lem:twoends}
Suppose the surviving maps over the two-ended component are power
cylinders with coarse contact profile \(\mu=(k_i)\) and the insertion
is pulled back from the main relative factor. For one prescribed
matching of the separately labelled ends the weight is \(\prod_i k_i\).
Against the main integral with both lists unlabelled it is
\(z_\mu=|\operatorname{Aut}\mu|\prod_i k_i\). This holds for target
index \(d=1\) and \(d=2\).
\end{lemma}
\begin{proof}
At one flag of coarse order \(k\) put \(a=(d,k)\), \(f=d/a\)
and \(c=k/a\). The primitive representable source index is \(f\);
in a uniformizer the contact has order \(c\). The inverse-band
matching diagonal of the rigidified evaluation sector
\(B(\mu_d/\mu_f)\) has degree \(a\): its fiber is the torsor of
isomorphisms between two given evaluations. Choose one auxiliary
twist \(R\) divisible by every \(c\). The target index becomes
\(dR\), the source index \(dR/k\), and the map of rigidified
sectors to the primitive sector is a \(\mu_c\)-gerbe. Comparison
of the diagonals contributes \(c\) per flag. The single target
gluing gerbe of degree \(1/R\) and the central multiplicity \(R\)
cancel on the common expansion base. Thus the flag factor is
\(ac=k\) \cite[\S\S5.5--5.10, Lemma~5.15 and Proposition~5.18]{af}.

With \(N=dR\), the cylinder is the index-\(k\) rotation quotient
\([{\PP}^1/\mu_{N/k}]\to[{\PP}^1/\mu_N]\). The simply connected
atlas gives a unique connected cover for this subgroup. Its log
differential is an isomorphism, and its deck group has order \(k\).
The two matching factors and its virtual degree give \(k^2/k=k\).
In particular for \(d=2\) and \(k\) even the two relative spin
clutches each carry weight \(k/2\): the deck action changes both
primitive signs together, and the two resulting orbits have
stabilizer of order \(k/2\). The pullback assumption gives the
same insertion on these two orbits.

Labelling equal cylinders introduces the divisor
\(A=|\operatorname{Aut}\mu|\) once, simultaneously on both ends.
The unlabelled main integral divides the separately labelled integral
by \(A^2\). Their ratio gives \(A\prod_i k_i\), as asserted.
\end{proof}

\subsection{The orientation at a contact}
In the primitive chart the sign line has an invariant matching fiber
at order \(k_i\) exactly when \(d\mid k_i\). Write \(I=\{i:d\mid k_i\}\).
On a surviving cylinder in \(I\), elimination of its constant at the
first end leaves the complex equation \(q_i-\sigma_i p_i=0\), where
\(p_i,q_i\) are the end values in square frames and \(\sigma_i=\pm1\).
Changing the target clutch negates each \(\sigma_i\).
On the cut main source Stokes' boundary inequality, after a common
phase rotation of the frames, is
\[
 \sum_{i\in I} k_i\operatorname{Re}(p_i^2-q_i^2)\ge0
\]
with equality only for the zero solution of the main equations.
For \(0\le u\le1\) impose instead
\[
 \operatorname{Re}p_i-u\sigma_i\operatorname{Re}q_i=0,\qquad
 \operatorname{Im}q_i-u\sigma_i\operatorname{Im}p_i=0 .
\]
Each term in the boundary sum then equals
\(k_i(u^2-1)((\operatorname{Re}q_i)^2+(\operatorname{Im}p_i)^2)\le0\).
The augmented operators remain injective throughout. Their endpoint
at \(u=0\) is independent of the clutch. At \(u=1\) the transformation
from the original complex row to the new two real rows has
determinant \(-\sigma_i\). Thus each changed invariant matching
fiber reverses the orientation once. For \(d=2\) the number of even
parts has parity \(m-r\); for \(d=1\) all \(r\) parts contribute.
The mixed insertion equals the repeated square multiplied by
\begin{equation}\label{eq:clutch}
 (-1)^{|I|}=\epsilon^m(-1)^{m-r}.
\end{equation}
Before this comparison one must keep the two spin matchings at an
even contact separate. After it the sign depends only on \(\mu\),
so finite transfer preserves the formula.

\subsection{The repeated square and the trace}
Let \(H\) denote the spin pulled back to the closed source after the
cylinders have been attached, and \(W_H=-R\pi_*H\). The square of
the Euler class of an oriented real \(2n\)-plane bundle is its top
Pontryagin class. Continuation to Dolbeault gives
\[
 e(V_H)^2=(-1)^n c_{2n}(W_H+\overline W_H).
\]
We compare this square with an ordinary index on the coarse source.
Choose a locally free spin \(L\) on the ordinary nodal target. Pull
it back to the twisted target. Its square agrees with that of \(H\)
in logarithmic frames, so \(F'=HL^{-1}\) on the target is a sign
line. There \([H]-[L]=[F']-1\) in the perfect \(K\)-group: express
\(L\) by divisors off the node and tensor the resulting point
sheaves by \(F'\). Pull this identity back along the universal map
before pushing it down to the parameter space. The zeroth
cohomologies of the sign and trivial lines have finite monodromy;
the positive even Chern characters of their first cohomologies
vanish by the Hodge relation on the coarse double torsor and its
invariant summand. Conjugation cancels the odd Chern characters.
Consequently
\(c_{2n}(W_H+\overline W_H)=c_{2n}(W_L+\overline W_L)\) rationally,
including on the twisted sectors. On normalizing the \(r\) power
cylinders the two complex rows at each cylinder cancel its single
constant, leaving \(W_L=W_B+\mathcal O^r\) after framing. Here \(W_B\)
is the ordinary relative index on the main coarse source. It
descends with the relative virtual class by
\cite[Theorem~2.1]{acw} and \cite[Theorem~1.1.2]{amw}.

For equivariant weight \(s_0\), the inverse antidiagonal pairing on
a profile of length \(r\) is \(z_\mu(-s_0^2)^r\);
the required degree of the equivariant Euler product of \(W_B\)
contributes \((-1)^{n-r}s_0^{-2r}c_{2n}(W_B+\overline W_B)\).
Their product is \(z_\mu\) times the repeated square. The factor
\(z_\mu\) agrees with Lemma~\ref{lem:twoends} on unlabelled relative
maps. This establishes the normalization and the orientation for the
full contraction, including both clutch types. The arithmetic genus
of the disconnected cut source, defined by \(1-\chi(\mathcal O)\),
is \(k-r\). The relative shift contributes \(2r\) to its genus
exponent, restoring \(2k-2\).

Let \(A_m\) be the balanced two-point operator of genus \(g-1\) and
levels \((b_0,b_0)\) in the antidiagonal relative theory. The
degree-\(m\) coefficient of \(Z_D\) is now
\(\epsilon^m\operatorname{Tr}(\Omega A_m)\), where
\(\Omega e_\mu=(-1)^{m-\ell(\mu)}e_\mu\) in the relative basis
indexed by contact profiles. In the idempotent indexed by
\(\lambda\), the coefficient of \(e_\mu\) is proportional to
\(\dim(\lambda)\chi^\lambda(\mu)\), with the remaining factor
depending only on \(m,\mu,s_0\). Tensoring the character with sign
therefore gives \(\Omega v_\lambda=v_{\lambda^\top}\), with unit
coefficient. The eigenvalue of \(A_m\) is
\(\prod_{u\in\lambda}X_{h(u)}^{b_0}\) by the handle and level
eigenvalues in \cite[\S7]{bp08}. Taking the trace proves the second
line of Theorem~\ref{thm:spin}.
